\documentclass[11pt]{article}

\usepackage[utf8]{inputenc}
\usepackage[T1]{fontenc}
\usepackage{geometry}
\usepackage{amsmath}
\usepackage{amssymb}
\usepackage{booktabs}
\usepackage{hyperref}
\usepackage{xurl}
\usepackage{xcolor}
\usepackage{listings}
\usepackage{enumitem}
\usepackage{parskip}
\usepackage{graphicx}
\graphicspath{{../images/}}

\geometry{margin=1in}

\lstset{
  language=Java,
  basicstyle=\ttfamily\small,
  keywordstyle=\color{blue},
  commentstyle=\color{gray},
  stringstyle=\color{teal},
  showstringspaces=false,
  columns=fullflexible,
  frame=single,
  framerule=0.2pt,
  breaklines=true
}

\setlist[itemize]{topsep=0.3em,itemsep=0.2em}
\setlist[enumerate]{topsep=0.3em,itemsep=0.2em}

% Simple per-section source line.
\newcommand{\notesources}[1]{\par\noindent{\small\color{gray}\textit{Sources:} \sloppy #1\par}}

% Unit-wise source strings for this subject (used by slides/notes).
% Path is relative to the lecture `latex/` folder (the compilation CWD).
% OOPS with Java (CSEG1044) — unit-wise sources
%
% These are short, student-facing reference strings intended to be used with:
%   - \slidesources{...}  (in beamer slides)
%   - \notesources{...}   (in lecture notes)
%
% Style inspiration: other subjects in My-Subjects/ use semicolon-separated sources
% and include an "accessed" date for web documentation.

\newcommand{\OOPSUnitISources}{Schildt, \textit{Java: A Beginner's Guide} (10e), McGraw-Hill (Java fundamentals + OOP basics).; Oracle Java Tutorials: Getting Started + Language Basics + Classes and Objects (accessed \today).; Java SE API docs (e.g., \texttt{String}, \texttt{Scanner}) (accessed \today).}

\newcommand{\OOPSUnitIISources}{Schildt, \textit{Java: The Complete Reference} (12e), McGraw-Hill (classes, inheritance, packages, interfaces).; Bloch, \textit{Effective Java} (3e), Addison-Wesley, 2018 (design choices; interfaces vs abstract classes).; Oracle Java Tutorials: Classes and Objects + Interfaces + Packages (accessed \today).; Java SE API docs (e.g., \texttt{Comparable}, \texttt{Runnable}) (accessed \today).}

\newcommand{\OOPSUnitIIISources}{Schildt, \textit{Java: The Complete Reference} (12e) (nested/inner classes, exceptions, threads, I/O).; Oracle Java Tutorials: Nested Classes + Exceptions + Concurrency + Basic I/O (accessed \today).; Java SE API docs (e.g., \texttt{Thread}, \texttt{AutoCloseable}) (accessed \today).}

\newcommand{\OOPSUnitIVSources}{Schildt, \textit{Java: The Complete Reference} (12e) (generics, lambdas, Swing, JDBC).; Bloch, \textit{Effective Java} (3e) (generics + lambdas).; Oracle Java Tutorials: Generics + Lambda Expressions + Swing + JDBC Basics (accessed \today).; Java SE API docs (e.g., \texttt{java.util.function}, \texttt{javax.swing}, \texttt{java.sql}) (accessed \today).}

\newcommand{\OOPSUnitVSources}{Schildt, \textit{Java: The Complete Reference} (12e) (Collections Framework + wrapper classes).; Oracle Java docs: Collections Framework overview (accessed \today).; Java SE API docs (\texttt{java.util} collections; wrapper classes in \texttt{java.lang}) (accessed \today).}

\newcommand{\OOPSUnitVISources}{Git SCM: \textit{Pro Git} (2e) (version control workflow), accessed \today.; GitHub Docs (repositories, issues, pull requests), accessed \today.; Oracle Java Tutorials: Swing + JDBC (accessed \today).; JUnit 5 User Guide (accessed \today).; Apache Maven Project: Getting Started (accessed \today).; Gradle User Manual: Getting Started (accessed \today).; UML class diagrams: UML 2.x overview/spec (accessed \today).}

\newcommand{\OOPSMiscWhyInterfacesSources}{Schildt, \textit{Java: The Complete Reference} (12e) (interfaces).; Bloch, \textit{Effective Java} (3e) (prefer interfaces where appropriate).; Oracle Java Tutorials: Interfaces (accessed \today).; Java SE API docs (e.g., \texttt{Comparable}, \texttt{Runnable}) (accessed \today).}



\title{Object Oriented Programming (CSEG1044)\\Unit 06 -- Lecture 05 Notes\\Swing UI Panels + Integration}
\author{Tofik Ali}
\date{\today}

\begin{document}
\maketitle
\tableofcontents

\section{Lecture Overview}
In Lecture-04 you built the database schema and a working DAO layer.
This lecture connects your Swing UI to your service/DAO so the application becomes a usable product.

\textbf{Core goal (Deliverable-4):} build one clear end-to-end flow:
\[
\text{UI} \rightarrow \text{Service (rules)} \rightarrow \text{DAO (JDBC)} \rightarrow \text{DB}
\rightarrow \text{UI refresh}
\]
\notesources{\OOPSUnitVISources}

\section{Syllabus Alignment (Unit 06)}
\textbf{Unit 06:} Swing UI, event handling, integration with service + JDBC persistence

\section{Learning Outcomes}
After this lecture, you should be able to:
\begin{itemize}
  \item Plan screens/panels and choose layout managers.
  \item Implement Swing event handlers that call service/DAO code cleanly.
  \item Validate input and show user-friendly messages.
  \item Complete Deliverable-4: working UI flow connected to persistence.
\end{itemize}

\section{Key Terms (Quick)}
\begin{itemize}
  \item \textbf{Panel}: a \texttt{JPanel} representing a screen or a section.
  \item \textbf{Layout manager}: positions components (\texttt{BorderLayout}, \texttt{GridBagLayout}, ...).
  \item \textbf{Event}: user action like click/typing; handled by listeners (e.g., \texttt{ActionListener}).
  \item \textbf{EDT}: Event Dispatch Thread; Swing paints UI and runs event handlers on this thread.
  \item \textbf{CardLayout}: layout used to switch between screens/panels.
  \item \textbf{Validation}: checks that input and business rules are correct before saving.
  \item \textbf{Layering}: UI $\rightarrow$ service $\rightarrow$ DAO $\rightarrow$ DB.
\end{itemize}

\section{Detailed Notes}
\subsection{Start with a screen/panel plan (before coding)}
\textbf{Step 1: list the screens/panels.} For a Resource Booking example:
\begin{itemize}
  \item Resource list (view/search resources)
  \item Add/Edit resource (short form)
  \item Booking form (create booking)
  \item Booking list (view bookings)
\end{itemize}

\textbf{Step 2: decide navigation.} Common patterns:
\begin{itemize}
  \item \texttt{CardLayout} for switching pages inside one main window.
  \item \texttt{JTabbedPane} for independent pages.
  \item \texttt{JDialog} for short add/edit forms.
\end{itemize}

\textbf{Step 3: pick one end-to-end use case first.} For Deliverable-4, pick one:
\begin{itemize}
  \item Add Resource $\rightarrow$ save in DB $\rightarrow$ refresh list
  \item Create Booking $\rightarrow$ save in DB $\rightarrow$ show booking in list
\end{itemize}

\subsection{Layout managers (practical cheat sheet)}
Avoid \texttt{null} layout.

\begin{tabular}{@{}p{0.23\linewidth}p{0.73\linewidth}@{}}
\toprule
\textbf{Layout} & \textbf{Good for} \\
\midrule
\texttt{BorderLayout} & Overall frame structure (menu/top bar/content). \\
\texttt{FlowLayout} & Small rows of buttons. \\
\texttt{GridLayout} & Uniform grids. \\
\texttt{BoxLayout} & Vertical/horizontal stacking; simple forms. \\
\texttt{GridBagLayout} & Flexible aligned forms (labels + fields). \\
\texttt{CardLayout} & Multi-screen apps (switch screens). \\
\bottomrule
\end{tabular}

\subsection{Event handling pattern (button click)}
Most flows follow this sequence:
\begin{enumerate}
  \item UI reads inputs.
  \item UI calls a service method.
  \item UI shows success/error message.
  \item UI refreshes the displayed data (table/list).
\end{enumerate}

\begin{lstlisting}
btnSave.addActionListener(e -> {
    try {
        service.createResource(txtName.getText(),
            txtType.getText(),
            Integer.parseInt(txtCapacity.getText()));
        JOptionPane.showMessageDialog(this, "Saved!");
        refreshTable();
    } catch (NumberFormatException ex) {
        JOptionPane.showMessageDialog(this, "Capacity must be a number.",
            "Validation", JOptionPane.WARNING_MESSAGE);
    } catch (ValidationException ex) {
        JOptionPane.showMessageDialog(this, ex.getMessage(),
            "Validation", JOptionPane.WARNING_MESSAGE);
    } catch (Exception ex) {
        JOptionPane.showMessageDialog(this, "Error: " + ex.getMessage(),
            "Error", JOptionPane.ERROR_MESSAGE);
    }
});
\end{lstlisting}

\textbf{Tip:} keep handlers short. Put rules in service methods so you can test them.

\subsection{Separation: UI vs service vs DAO (what goes where?)}
\textbf{UI layer:}
\begin{itemize}
  \item reads UI components and converts strings to types,
  \item shows messages and updates tables/lists,
  \item calls service methods.
\end{itemize}

\textbf{Service layer:}
\begin{itemize}
  \item validates inputs and business rules,
  \item calls DAO methods,
  \item throws clear exceptions such as \texttt{ValidationException}.
\end{itemize}

\textbf{DAO layer:}
\begin{itemize}
  \item performs SQL CRUD and mapping only,
  \item uses \texttt{PreparedStatement} and \texttt{try-with-resources}.
\end{itemize}

\subsection{Validation + user-friendly error messages}
Validate at two levels:
\begin{itemize}
  \item \textbf{UI validation (fast feedback)}: empty fields, number parsing.
  \item \textbf{Service validation (complete)}: all rules, including cross-field rules.
\end{itemize}

\textbf{Example rules:}
\begin{itemize}
  \item Resource name required (3--100 chars)
  \item Capacity must be $>$ 0
  \item Booking start time must be before end time
  \item No overlapping bookings for the same resource (service rule)
\end{itemize}

\subsection{EDT rule + long tasks}
Swing is single-threaded:
\begin{itemize}
  \item UI painting and event handlers run on the \textbf{EDT}.
  \item If you call the DB on the EDT and it is slow, the UI freezes.
\end{itemize}

If DB operations are slow, use \texttt{SwingWorker}:
\begin{lstlisting}
new SwingWorker<List<Resource>, Void>() {
    protected List<Resource> doInBackground() throws Exception {
        return service.listResources();
    }
    protected void done() {
        try { updateTable(get()); }
        catch (Exception ex) { showError(ex); }
    }
}.execute();
\end{lstlisting}

\notesources{\OOPSUnitVISources}

\section{Worked Example: Add Resource $\rightarrow$ DB $\rightarrow$ refresh}
This worked example demonstrates the smallest complete integration.
\textbf{Imports are omitted} for readability.

\subsection{Service (validation + DAO call)}
\begin{lstlisting}
public class ValidationException extends RuntimeException {
    public ValidationException(String msg) { super(msg); }
}

public class ResourceService {
    private final ResourceDao dao;
    public ResourceService(ResourceDao dao) { this.dao = dao; }

    public void createResource(String name, String type, int capacity) throws Exception {
        if (name == null || name.trim().isEmpty()) throw new ValidationException("Name is required.");
        if (type == null || type.trim().isEmpty()) throw new ValidationException("Type is required.");
        if (capacity <= 0) throw new ValidationException("Capacity must be > 0.");
        dao.insert(new Resource(0, name.trim(), type.trim(), capacity));
    }

    public List<Resource> listResources() throws Exception {
        return dao.findAll();
    }
}
\end{lstlisting}

\subsection{List panel (JTable + refresh)}
\begin{lstlisting}
public class ResourceListPanel extends JPanel {
    private final ResourceService service;
    private final DefaultTableModel model =
        new DefaultTableModel(new Object[]{"ID","Name","Type","Cap"}, 0);

    public ResourceListPanel(ResourceService service) {
        this.service = service;
        setLayout(new BorderLayout());

        JTable table = new JTable(model);
        add(new JScrollPane(table), BorderLayout.CENTER);

        JButton btnAdd = new JButton("Add Resource");
        btnAdd.addActionListener(e -> openAddDialog());
        add(btnAdd, BorderLayout.SOUTH);

        refreshTable();
    }

    private void refreshTable() {
        try {
            model.setRowCount(0);
            for (Resource r : service.listResources()) {
                model.addRow(new Object[]{r.getId(), r.getName(), r.getType(), r.getCapacity()});
            }
        } catch (Exception ex) {
            JOptionPane.showMessageDialog(this, ex.getMessage(), "Error", JOptionPane.ERROR_MESSAGE);
        }
    }

    private void openAddDialog() {
        AddResourceDialog dlg = new AddResourceDialog(service, this::refreshTable);
        dlg.setVisible(true);
    }
}
\end{lstlisting}

\subsection{Add dialog (short form)}
\begin{lstlisting}
public class AddResourceDialog extends JDialog {
    public AddResourceDialog(ResourceService service, Runnable onSaved) {
        setTitle("Add Resource");
        setModal(true);
        setSize(420, 240);
        setLocationRelativeTo(null);

        JTextField txtName = new JTextField(20);
        JTextField txtType = new JTextField(20);
        JTextField txtCap = new JTextField(6);

        JButton btnSave = new JButton("Save");
        btnSave.addActionListener(e -> {
            try {
                service.createResource(txtName.getText(), txtType.getText(),
                    Integer.parseInt(txtCap.getText()));
                JOptionPane.showMessageDialog(this, "Saved!");
                onSaved.run();
                dispose();
            } catch (NumberFormatException ex) {
                JOptionPane.showMessageDialog(this, "Capacity must be a number.",
                    "Validation", JOptionPane.WARNING_MESSAGE);
            } catch (ValidationException ex) {
                JOptionPane.showMessageDialog(this, ex.getMessage(),
                    "Validation", JOptionPane.WARNING_MESSAGE);
            } catch (Exception ex) {
                JOptionPane.showMessageDialog(this, ex.getMessage(),
                    "Error", JOptionPane.ERROR_MESSAGE);
            }
        });

        JPanel form = new JPanel(new GridLayout(3, 2, 8, 8));
        form.add(new JLabel("Name")); form.add(txtName);
        form.add(new JLabel("Type")); form.add(txtType);
        form.add(new JLabel("Capacity")); form.add(txtCap);

        add(form, BorderLayout.CENTER);
        add(btnSave, BorderLayout.SOUTH);
    }
}
\end{lstlisting}

\section{Deliverable-4 checklist}
\begin{itemize}
  \item At least 2 UI screens/panels (list + form) with navigation.
  \item One complete end-to-end flow: UI $\rightarrow$ service $\rightarrow$ DAO $\rightarrow$ DB.
  \item Validations + user-friendly messages (no stack traces for users).
  \item Evidence in GitHub: issues, commits, PR; short README ``how to run''.
\end{itemize}

\section{In-Class Exercises (with brief solutions)}
\begin{enumerate}
  \item Which use case will you implement first for end-to-end proof?\par
  \textbf{Answer:} A create flow that also lets you see data back (create + list).

  \item Where should the rule ``booking end time must be after start time'' live?\par
  \textbf{Answer:} In service (business rule). UI may do basic date checks too.

  \item Your UI freezes on ``Refresh''. What is the cause and fix?\par
  \textbf{Answer:} DB call on EDT; use \texttt{SwingWorker} or background thread.

  \item Should DAO exceptions be shown raw to the user?\par
  \textbf{Answer:} No; show a friendly message and keep details for debugging/logs.
\end{enumerate}

\section{Checkpoint Questions (with sample answers)}
\textbf{Q1.} Why separate UI/service/DAO?\par
\textbf{Sample answer:} It keeps code maintainable and testable; each layer has one responsibility.

\vspace{0.6em}
\textbf{Q2.} What is the EDT and why does it matter?\par
\textbf{Sample answer:} Swing uses the EDT for painting and events; long work on EDT freezes the UI.

\end{document}
